\section{邻居信息参与价值更新的动机与机制}

空间公共物品博弈中，个体的收益同时受到自身行为和周围个体行为的影响。一次合作之后获得较低收益，并不能单独说明合作在后续交互中仍然不利。邻居的行为及其获得的反馈，为个体判断当前策略提供了额外的参照。本研究据此在 Q-learning 的价值更新中引入邻居影响，考察相对表现信息如何改变个体的学习过程，以及这种改变能否促进群体合作。

模型将博弈交互范围与信息感知范围分别设置。每个智能体参与以自身及四个一阶邻居为中心的五个公共物品博弈组；感知参数 $M$ 则决定可用于状态构建和邻居评价的信息范围。当 $M=2$ 时，感知集合包含自身和十二个邻居，其中可作为参考对象的邻居有十二个。扩大感知范围改变了可获取的信息，并未扩大每个博弈组的成员范围。

个体维护自己的 Q 表，并依据感知范围内的平均声誉构造状态。动作执行后，模型将归一化的博弈收益与声誉变化加权组合，得到用于学习的综合奖励。随后先进行标准的时序差分更新，再计算邻居影响。因此，邻居比较所使用的是综合奖励，不能直接等同于未经处理的博弈收益。

邻居影响机制从候选邻居中选出综合奖励最高的个体。只有该邻居的奖励高于自身时，附加修正才生效：双方动作一致时，增加当前状态下已执行动作的 Q 值；动作不一致时，减小该 Q 值。修正幅度由正奖励差和控制参数 $\kappa$ 共同决定，并使用全局最大邻居奖励差进行归一化。完成更新后，智能体仍依据自己的 Q 表，通过 $\epsilon$-greedy 策略选择后续动作。邻居信息由此影响个体对动作价值的判断，而动作选择仍由个体自身的策略完成。

上述规则以邻居报告能够反映其真实表现为前提。若报告值被人为抬高，按最大综合奖励选择参考对象便可能反复选中异常发送者，进而影响价值更新。这一局限引出了后续研究：在保留个体 Q-learning 决策过程的基础上，进一步研究邻居消息的选择方式及其修正强度。该问题能否得到缓解，需要在明确的消息异常模型下，通过对照实验检验。
